  \def\stat{gr-nik}





\def\tit{КОМПЛЕКСНОЕ ОБЕСПЕЧЕНИЕ ИНФОРМАЦИОННОЙ БЕЗОПАСНОСТИ ЧАСТНЫХ 


ОБЛАЧНЫХ ВЫЧИСЛИТЕЛЬНЫХ~СРЕД$^*$}





\def\titkol{Комплексное обеспечение ИБ %информационной безопасности 


частных  облачных вычислительных сред}








\def\aut{А.\,А.~Грушо$^1$, А.\,В.~Николаев$^2$, В.\,О.~Писковский$^3$, 


В.\,В.~Сенчило$^4$, Е.\,Е.~Тимонина$^5$}





\def\autkol{А.\,А.~Грушо, А.\,В.~Николаев, В.\,О.~Писковский и~др.}


%$^3$, В.\,В.~Сенчило$^4$, Е.\,Е.~Тимонина$^5$}





\titel{\tit}{\aut}{\autkol}{\titkol}





\index{Грушо А.\,А.}


\index{Николаев А.\,В.}


\index{Писковский В.\,О.} 


\index{Сенчило В.\,В.}


\index{Тимонина Е.\,Е.}


\index{Grusho A.\,A.}


\index{Nikolaev A.\,V.}


\index{Piskovski V.\,O.}


\index{Senchilo V.\,V.} 


\index{Timonina E.\,E.}








{\renewcommand{\thefootnote}{\fnsymbol{footnote}}


\footnotetext[1]


{Работа частично поддержана РФФИ (проект 18-29-03081).} 








\renewcommand{\thefootnote}{\arabic{footnote}}


\footnotetext[1]{Институт проблем информатики Федерального исследовательского 


центра <<Информатика и~управление>> 


Российской академии наук, \mbox{grusho@yandex.ru}}


\footnotetext[2]{Федеральный исследовательский центр химической физики им.\ Н.\,Н.~Семёнова Российской академии наук, 


\mbox{gentoorion@mail.ru}}


\footnotetext[3]{Институт проблем информатики Федерального исследовательского 


центра <<Информатика и~управление>> 


Российской академии наук, \mbox{vpvp80@yandex.ru}}


\footnotetext[4]{Институт проблем информатики Федерального исследовательского 


центра <<Информатика и~управление>> 


Российской академии наук, \mbox{volodias@mail.ru}}


\footnotetext[5]{Институт проблем информатики Федерального исследовательского 


центра <<Информатика и~управление>> 


Российской академии наук, \mbox{eltimon@yandex.ru}}





 \label{st\stat}





 \thispagestyle{headings}


 


 \vspace*{-12pt}  


  





  


  \Abst{Рассматривается задача построения безопасного корпоративного облака, в~котором 


используются небезопасные компоненты. Наиболее важным представляется решение для 


удаленной работы сотрудников организации. Решение основано на выделении простого 


подоблака, реализующего безопасные взаимодействия удаленных сотрудников между собой, с~облачными сервисами и~ресурсами глобальной сети Интернет.}


  


  \KWE{информационная безопасность; облачные вычисления; безопасность удаленной 


работы на небезопасных терминалах}





\DOI{10.14357\!/\!08696527200407}  





\vskip 10pt plus 9pt minus 6pt





 \thispagestyle{headings}


 


 \vspace*{-24pt}


  





  \section{Введение }


  


  Различные организации все шире используют облачные вычисления. Вынос 


точки доступа в~частное облако за пределы контролируемой зоны (КЗ) всегда 


связан с~целым рядом рисков информационной безопасности (ИБ). В~качестве 


такой точки доступа может быть использован ноутбук, персональный компьютер, 


планшет или телефон. Удаленное рабочее место может содержать деловую 


переписку, важные документы, статистику. Его легко потерять, лишиться при 


транспортировке и~пр. Таким образом, обработка важной или секретной 


информации на таком рабочем месте вне КЗ проблематична. Кроме 


несанкционированного доступа к~конфиденциальной информации, хранимой на 


рабочей станции, подключение ее к~домашней сети может завершиться взломом 


и~инфицированием с~последующим неконтролируемым проникновением 


в~офисную информационную среду или инфраструктуру организации. Без 


профессиональных средств защиты обнаружить подобную атаку практически 


невозможно, как и~невозможно предотвратить использование сторонних средств 


хранения и~обмена информации, ин\-тер\-нет-ре\-сур\-сов. Служба 


ИБ компании обычно не имеет возможности проконтролировать 


удаленные рабочие места, что делает такие сценарии реализации угроз еще более 


критичными с~точки зрения последствий. Контроль действий при обмене 


информацией становится еще более сложной задачей. Очевидно, что для 


обеспечения удаленного рабочего места, безопасного с~точки зрения интересов 


компании, нужен новый подход к~организации эффективной системы защиты, 


которая создавалась бы с~учетом вышеперечисленных угроз. При этом система 


должна быть простой в~эксплуатации, а лучше~--- скрытой, управление не должно 


требовать высокой квалификации обслуживающего персонала.


  


  \section{Архитектура защиты }


  


  В течение ряда лет в~Российской Федерации проводились исследования по 


разработке архитектуры защищенной информационной системы (ИС) из 


не\-безопас\-ных компонентов~[1, 2]. Ряд этих результатов можно применить 


и~к~построению корпоративного облака. В работе предложено в~отдельное 


подоблако вынести: 


  \begin{itemize}


\item удаленное взаимодействие сотрудников корпорации между собой;


\item удаленное использование ресурсов компании; 


\item взаимодействие с~ресурсами глобальной сети. 


\end{itemize}





  Такое подоблако позволит локализовать и~изолировать опасные функции, 


связанные с~взаимодействиями и~сетевыми обменами. Отметим, что остальная 


часть облачных сервисов также окажется в~изоляции, что позволит использовать 


в~облаке небезопасные компоненты. Идея максимальной изоляции подсистем, 


построенных из небезопасных элементов, стала одной из основных в~теории 


построения защищенных систем из незащищенных элементов~[3, 4].


  


  Этот подход частично пересекается с~новой разработкой американских ученых, 


а~именно: с~концепцией в~области кибербезопасности~--- <<нулевое доверие>> 


(Zero Trust, или ZT)~[5]. Этот подход предполагает, что каждая сеть взломана, 


каждая машина скомпрометирована и~каждый пользователь (вольно или невольно) 


становится либо нарушителем, либо объектом интереса нарушителя. Никому 


и~ничему в~сети с~нулевым доверием нельзя доверять, пока не доказано обратное. 


  


  Национальный институт стандартов и~технологий США (NIST)~[5] определяет 


ZT как парадигму ки\-бер\-безопас\-ности, направленную на защиту ресурсов 


в~предположении, что уровень доверия никогда не предоставляется по 


умолчанию, а~постоянно оценивается. Архитектура сети, соответствующая ZT, 


представляет собой план кибербезопасности предприятия, который использует 


концепции нулевого доверия и~охватывает отношения компонентов, планирование 


рабочего процесса и~политики доступа. Таким образом, ZT-безопас\-ность~--- это 


совокупность взаимосвязанных политик, практик, программного обеспечения 


(ПО) и~оборудования.


  


  ZT не выступает альтернативой традиционным методам безопасности, 


а~является развитием хорошо известных и~зарекомендовавших себя технологий. 


Повысить уровень доверия к~среде, в~которой работает пользователь, ZT 


предлагает с~помощью ряда известных мер, в~числе которых: 


  \begin{itemize}


\item установка и~поддержка антивирусного ПО; 


\item установка и~корректная настройка межсетевого экрана; 


\item корректная настройка процессов обновления системного и~прикладного 


ПО;


\item установка аппаратного модуля доверенной загрузки. 


\end{itemize}





  Для создания сети в~рассматриваемых предположениях необходимо наличие 


ряда автоматизированных служб и~компонентов поддержки безопасности. Все эти 


элементы служат инструментами или процедурами, которые уже успешно 


существуют, эксплуатируются и~повторяют концепцию  


про\-грам\-мно-ком\-му\-ти\-ру\-емой сети (ПКС). Так, компонент принятия 


решения о политике (PDP~--- Policy Decision Point) решает, безопасен ли 


компьютер или веб-тра\-фик, и~после этого соответственно предоставляет или 


отменяет доступ. Механизм политик безопас\-ности (ПБ)~[6] использует все 


доступные, в~том числе и~внешние, источники данных, применимые 


к~потребностям организации для определения безопасности на основе ПБ. Так, 


решение о безопасности сетевого компонента принимается на основе 


информации, предоставляемой системами. Система непрерывной диагностики 


и~компенсации последствий считается главной. Система собирает информацию 


о~текущем состоянии сетевых компонентов и~применяет обновления 


к~конфигурации и~программным компонентам, предоставляет информацию об 


элементе, отправляющем запрос на доступ, например работает ли он 


с~соответствующей исправленной операционной системой и~приложениями или 


есть ли у~актива ка\-кие-ли\-бо известные уязвимости.


  


   Основным недостатком описанного подхода является наличие сложной 


интеллектуальной системы отслеживания и~динамического принятия решений. 


В~рамках достижений информационных технологий в~РФ такую систему создать 


сложно. Второй недостаток американского подхода состоит в~том, что он 


опирается на стандартный набор механизмов безопасности, которые не всегда 


дают надежность защиты. В~отечественных разработках предложена концепции 


контроля соединений распределенной системы на основе метаданных, которая 


позволяет обеспечить более высокий уровень ИБ~[7].


  


  \section{Безопасные взаимодействия }


  


  В качестве примера применения частного подоблака для организации 


безопасных взаимодействий рассмотрим реализацию частного облака, 


представленного на рисунке. 





\begin{figure}


\vspace*{1pt}


\begin{center}


\mbox{%


\epsfxsize=107mm %8.78mm  


\epsfbox{gru-1.eps}


}


\end{center}


%\vspace*{-9pt}





{\small Архитектура защищенного корпоративного облака с~по\-мощью безопас\-но\-го 


коммуникационного подоблака}


\vspace*{-9pt}


\end{figure}





  В качестве пользовательского терминала используется виртуальная машина 


(или контейнер), на которой реализован клиент для терминальных служб, 


развернутых в~частном облаке. Привлекает идея загрузки в~недоверенный 


терминал пользователя <<чистой>> рабочей среды с~внешнего носителя или 


облачного ресурса. Будем называть устройство безопасного доступа 


пользовательского терминала к~облачным ресурсам \textit{абонентским облачным 


интерфейсом} (АОИ). Рабочая станция (РС) пользователя не имеет возможности 


непосредственного выхода в~глобальную сеть. Единственным каналом служит 


защищенный канал из АОИ, соединяющий РС1 с~частным подоблаком 


и~реализующий тотальное шифрование информации, передаваемой по этому 


каналу на уникальном для АОИ ключе. 


  


  В соответствии с~вышесказанным, РС1 с~установленным АОИ запрашивает 


нужное взаимодействие в~подоблаке контроля взаимодействий. Между 


коммутатором, служащим точкой доступа в~облако, и~АОИ устанавливается 


защищенное соединение на ключе, используемом только между этой 


РС и~шифратором на коммутаторе. Каждый легальный АОИ имеет 


уникальный симметричный ключ для связи с~коммутатором. Успешное 


расшифрование обращения к~коммутатору служит надежной аутентификацией 


АОИ и~инициирует продолжение организации взаимодействия.


  


  Если взаимодействие предполагает связь с~другим сотрудником корпорации, то 


осуществляется его вызов через аналогичный терминал этого сотрудника на его 


симметричном ключе. Для их взаимодействия потребуется перешифрование 


обмена информацией с~одного ключа на другой. Описанная схема взаимодействия~--- 


простейшая, она оставляет возможность контроля взаимодействия. Для 


маршрутизации трафика внутри облака коммутаторы строят индивидуальные для 


каждого узла защищенные соединения с~помощью развернутой инфраструктуры 


средств шифрования предприятия.


 


  \section{Заключение }


  


  Построение современных корпоративных систем с~защищенными терминалами 


удаленных друг от друга пользователей может опираться на незащищенные 


РС и~мобильные устройства, в~которые встроен единственный 


доверенный блок, реализующий интерфейс небезопасного компьютера без 


воз\-мож\-ности выхода в~сеть и~сетевого интерфейса с~единственным адресом 


взаимодействия. Этот адрес предполагает связь с~примитивным подоблаком 


корпоративной сис\-те\-мы, реализующей коммутатор для организации 


взаимодействия удаленных пользователей друг с~другом, с~сервисами 


корпоративного облака и~ресурсами внешней глобальной сети. В~за\-ви\-си\-мости от 


ценностей активов организации может использоваться гарантированное  или легкое шифрование.


  


{\small\frenchspacing


{%\baselineskip=10.8pt


\addcontentsline{toc}{section}{Литература}


\begin{thebibliography}{9} 


\bibitem{1-gn}


\Au{Grusho A.\,A.} Composition of the trusted computer security using untrusted elements~/\!\!/ 


Probabilistic methods in discrete mathematics~/ Eds. V.\,F.~Kolchin, 


V.\,Ya.~Kozlov, Yu.\,L.~Pavlov, Yu.\,V.~Prokhorov.~--- Utrecht, 


Netherlands: VSP, 1997. P.~231--236. doi: 10.1515/9783112314074-022.


\bibitem{2-gn}


\Au{Тимонина Е.\,Е.} Анализ угроз скрытых каналов и~методы построения гарантированно 


защищенных распределенных автоматизированных систем: Дис.\ \ldots\ докт. техн. наук.~--- М., 2004. 


204~с.


\bibitem{3-gn}


\Au{Грушо А.\,А., Володин А.\,В., Тимонина~Е.\,Е.} Безопасный интерфейс с~глобальной сетью из 


ненадежных в~смысле безопасности элементов~/\!\!/ Информационные технологии в~науке, 


образовании, телекоммуникации, бизнесе (весенняя сессия): Тр. Междунар. конф.~--- 


Запорожье, Украина: Лаборатория издательских технологий и~компьютерной графики 


Запорожского государственного университета, 2001. С.~149--151.


\bibitem{4-gn}


\Au{Грушо А.\,А., Тимонина Е.\,Е.} Гарантированно защищенные базы данных, построенные на 


недоверенных с~точки зрения безопасности элементах~/\!\!/ Проб\-ле\-мы безопас\-ности 


и~противодействия терроризму: Мат-лы конф.~--- М.: МЦНМО, 2006. С.~335--348.


\bibitem{5-gn}


\Au{Borchert O., Mitchell S., Connelly~S., Rose~S.} Zero trust architecture~/\!\!/ Draft (2nd) NIST 


Special Publication 800-207, 2000. {\sf https://csrc.nist.gov/publications/\linebreak detail/sp/800-207/final}. 


\bibitem{6-gn}


\Au{Грушо А.\,А., Применко Э.\,А., Тимонина~Е.\,Е.} Теоретические основы компьютерной 


безопас\-ности.~--- М.: Академия, 2009. 272~с.


\bibitem{7-gn}


\Au{Grusho A.\,A., Timonina E.\,E., Shorgin~S.\,Ya}. Modelling for ensuring information security of 


the distributed information systems~/\!\!/ 31th European Conference on Modelling and Simulation 


Proceedings.~--- Dudweiler, Germany: Digitaldruck Pirrot GmbH, 2017. P.~656--660.


{\sf 


 http://www.scs-europe.net/dlib/2017/ecms2017acceptedpapers/0656-probstatECMS20170026.pdf}.


\end{thebibliography}


} }








\vspace*{-6pt}





\hfill{\small\textit{Поступила в~редакцию 12.09.20}}














%\vspace*{6pt}





%\hrule





%\vspace*{2pt}





\newpage





\vspace*{-28pt}





%\hrule





%\vspace*{8pt}





\def\tit{END-TO-END INFORMATION SECURITY OF~PRIVATE CLOUD COMPUTING}





\def\titkol{End-to-end information security of~private cloud computing}











\def\aut{A.\,A.~Grusho$^1$, A.\,V.~Nikolaev$^2$, V.\,O.~Piskovski$^1$, V.\,V.~Senchilo$^1$, 


and~E.\,E.~Timonina$^1$}





\def\autkol{A.\,A.~Grusho, A.\,V.~Nikolaev, V.\,O.~Piskovski, et al.}


%V.\,V.~Senchilo$^1$,  and~E.\,E.~Timonina$^1$}





\titel{\tit}{\aut}{\autkol}{\titkol}





\vspace*{-11pt}








\noindent


$^1$Institute of Informatics Problems, Federal Research Center ``Computer Science\linebreak


$\hphantom{^1}$and Control'' of 


the Russian Academy of Sciences; 44-2~Vavilov Str., Moscow\linebreak


$\hphantom{^1}$119133, Russian Federation








\noindent


$^2$N.\,N.~Semenov Federal Research Center for Chemical Physics of the Russian\linebreak


$\hphantom{^1}$Academy of Sciences,


4~Kosygin  Str., Moscow 119991, Russian Federation





\def\leftfootline{\small{\textbf{\thepage}


\hfill Sistemy i Sredstva Informatiki~---


Systems and Means of Informatics 2020 vol~30 no\,4}


}%


 \def\rightfootline{\small{Sistemy i Sredstva Informatiki~---


 Systems and Means of Informatics 2020 vol~30 no\,4


\hfill \textbf{\thepage}}}


  


  


\vspace*{3pt}


 








\Abste{The paper discusses building of a secure enterprise cloud that uses unsecure components. The 


most important solution is for the organization's employees to work remotely. The solution is based on 


provision of a simple subcloud that implements secure interactions of remote employees among 


themselves, cloud services, and interaction with resources in the global Internet.}





\KWE{information security; cloud computing; security of remote work on unsecure terminals}











\DOI{10.14357\!/\!08696527200407}  





\vspace*{-12pt}








\Ack


\noindent


The paper was partially supported by RFBR, project  


No.\,18-29-03081.


 





%\vspace*{-12pt}





\renewcommand{\bibname}{\protect\rmfamily References}


%\renewcommand{\bibname}{\large\protect\rm References}





{\small\frenchspacing


{%\baselineskip=10.8pt


\addcontentsline{toc}{section}{References}


\begin{thebibliography}{9}


\bibitem{1-gn-1}


\Aue{Grusho, A.\,A.} 1997. Composition of a trusted computer security system on the base of 


untrusted elements. \textit{Probabilistic methods in discrete mathematics}. Eds. V.\,F.~Kolchin, 


V.\,Ya.~Kozlov, Yu.\,L.~Pavlov, and Yu.\,V.~Prokhorov. Utrecht, Netherlands: VSP. 231--236.  doi: 10.1515/9783112314074-022.


\bibitem{2-gn-1}


\Aue{Timonina, E.\,E.} 2004.  Analiz ugroz skrytykh kanalov i~metody postroeniya garantirovanno 


zashchishchennykh raspredelennykh avtomatizirovannykh sistem [The analysis of threats of covert 


channels and methods of creation of guaranteed protected distributed automated systems]. Moscow: 


Russian State University for the Humanities. D.Sc. Diss. 204~p.


\bibitem{3-gn-1}


\Aue{Grusho, A.\,A., A.\,V. Volodin, and E.\,E.~Timonina.} 2001. Bezopasnyy interfeys s~glo\-bal'\-noy 


set'yu iz nenadezhnykh v~smysle bezopasnosti elementov [Secure interface with WAN from untrusted 


elements]. \textit{Informatsionnye tekhnologii v~nauke, obrazovanii, 


telekommunikatsii, biznese (vesennyaya sessiya): Tr. Mezhdunar. konf.} [Conference (International)  ``Information 


Technologies in Science, Education, Telecommunications, Business (spring session)'' Proceedings]. 


Zaporozhye: Laboratory of Publishing Technologies and Computer Graphics of Zaporozhye State University. 


149--151.


\bibitem{4-gn-1}


\Aue{Grusho, A.\,A., and E.\,E.~Timonina.} 2006. Garantirovanno zashchishchennye bazy dannykh, 


postroennye na nedoverennykh s~tochki zreniya bezopasnosti elementakh [Guaranteed secure 


databases built on insecure elements]. \textit{Problemy bezopas\-nosti i~protivodeystviya 


terrorizmu: Mat-ly konf.} [Conference ``Security and Counter-Terrorism Issues'' Proceedings]. 


Moscow: MCCME. 335--348.


\bibitem{5-gn-1}


\Aue{Borchert, O., S.~Mitchell, S.~Connelly, and S.~Rose.} 2000. SP 800-207: Zero trust 


architecture. {NIST}. Available at: {\sf https://csrc.nist.gov/publications/detail/sp/800-207/final} 


(accessed September~12, 2020).


\bibitem{6-gn-1}


\Aue{Grusho, A., E. Primenko, and E.~Timonina.} 2009. \textit{Teoreticheskie osnovy komp'yuternoy 


bezopasnosti} [Theoretical bases of computer security]. Moscow: Academy. 272~р.


\bibitem{7-gn-1}


\Aue{Grusho, A.\,A., E.\,E.~Timonina, and S.\,Ya.~Shorgin.} 2017. Modelling for ensuring 


information security of the distributed information systems. \textit{31th European Conference on 


Modelling and Simulation Proceedings}. Dudweiler, Germany: Digitaldruck Pirrot GmbH. 656--660. 


Available at: {\sf 


 http://www.scs-europe.net/dlib/2017/ecms2017acceptedpapers/0656-probstatECMS20170026.pdf} 


(accessed September12, 2020).


\end{thebibliography}


} }





%\vspace*{-3pt}





\hfill{\small\textit{Received September 12, 2020}} 





%\vspace*{-16pt}





\Contr





\noindent


\textbf{Grusho Alexander A.} (b.\ 1946)~--- Doctor of Science in physics and mathematics, professor, 


principal scientist, Institute of Informatics Problems, Federal Research Center ``Computer Science 


and Control'' of the Russian Academy of Sciences; 44-2~Vavilov Str., Moscow 119133, Russian 


Federation; \mbox{grusho@yandex.ru}





\vspace*{3pt}





\noindent


\textbf{Nikolaev Andrei V.} (b.\ 1973)~--- Candidate of Science (PhD) in physics and mathematics, 


scientist, N.\,N.~Semenov Federal Research Center for Chemical Physics of the Russian Academy of Sciences, 


4~Kosygin Str., Moscow 119991, Russian Federation; \mbox{gentoorion@mail.ru}





\vspace*{3pt}





\noindent


\textbf{Piskovski Viktor O.} (b.\ 1963)~--- Candidate of Science (PhD) in physics and mathematics, 


senior scientist, Institute of Informatics Problems, Federal Research Center ``Computer Science and 


Control'' of the Russian Academy of Sciences; 44-2~Vavilov Str., Moscow 119133, Russian 


Federation; \mbox{vpvp80@yandex.ru}





\vspace*{3pt}





\noindent


\textbf{Senchilo Vladimir V.} (b.\ 1963)~--- scientist, Institute of Informatics Problems, Federal 


Research Center ``Computer Science and Control'' of the Russian Academy of Sciences; 


 44-2~Vavilov Str., Moscow 119133, Russian Federation; \mbox{volodias@mail.ru} 





\vspace*{3pt}





\noindent


\textbf{Timonina Elena E.} (b.\ 1952)~--- Doctor of Science in technology, professor, leading scientist, 


Institute of Informatics Problems, Federal Research Center ``Computer Science and Control'' of the 


Russian Academy of Sciences; 44-2~Vavilov Str., Moscow 119133, Russian Federation; 


\mbox{eltimon@yandex.ru}





\label{end\stat}





\renewcommand{\bibname}{\protect\rm Литература}


